\chapter{Ważne grupy macierzowe i ich algebry Liego}
\label{app:grupy-algebry-macierzowe}

\index{grupa Liego!macierzowa}
\index{algebra Liego!macierzowa}
\index{komutator macierzy}
\index{nawias Liego}

W dodatku zbieramy najważniejsze klasyczne grupy macierzowe oraz
odpowiadające im algebry Liego. Pozwala to uniknąć oddzielnego
omawiania tych samych obiektów w dwóch miejscach.

Jeśli $G$ jest macierzową grupą Liego, jej algebra Liego jest
przestrzenią styczną w elemencie neutralnym:
\[
\mathfrak g=T_IG.
\]
Dla grup macierzowych można ją traktować jako odpowiednią przestrzeń
macierzy z nawiasem
\[
[X,Y]:=XY-YX.
\]
W szczególności
\[
\dim G=\dim\mathfrak g.
\]

W przypadku grup zespolonych będziemy zaznaczać, czy wymiar liczony
jest nad $\mathbb C$, czy nad $\mathbb R$. Jeśli $G$ jest zespoloną
grupą Liego, to
\[
\dim_{\mathbb R}G
=
2\dim_{\mathbb C}G.
\]
Dalej $n\geq1$, $\mathbb K\in\{\R,\C\}$ oraz
$\mathbb K_*=\mathbb K\setminus\{0\}$.
Wymiar ilorazu przez domkniętą podgrupę liczymy
według Twierdzenia~\ref{thm:klein-iloraz-liego}.


% ============================================================
% GL
% ============================================================

\section{Ogólna grupa liniowa i algebra
\texorpdfstring{$\mathfrak{gl}(n,\mathbb K)$}{gl(n,K)}}
\label{sec:app-grupa-gl}

\index{grupa Liego!GL@$GL(n,\mathbb K)$}
\index{ogólna grupa liniowa}
\index{algebra Liego!gl@$\mathfrak{gl}(n,\mathbb K)$, ogólna liniowa}

Niech $\mathbb K=\mathbb R$ lub $\mathbb C$.

\begin{definicja}
\emph{Ogólną grupą liniową} nazywamy
\[
GL(n,\mathbb K)
=
\{A\in M_n(\mathbb K):\det A\neq0\}.
\]
Jej algebra Liego jest równa
\[
\mathfrak{gl}(n,\mathbb K)
=
M_n(\mathbb K).
\]
\end{definicja}

Ponieważ $GL(n,\mathbb K)$ jest otwartym podzbiorem
$M_n(\mathbb K)$, otrzymujemy
\[
\dim_{\mathbb R}GL(n,\mathbb R)
=
\dim_{\mathbb R}\mathfrak{gl}(n,\mathbb R)
=
n^2,
\]
natomiast
\[
\dim_{\mathbb C}GL(n,\mathbb C)
=
\dim_{\mathbb C}\mathfrak{gl}(n,\mathbb C)
=
n^2.
\]
Jako rzeczywiste rozmaitości
\[
\dim_{\mathbb R}GL(n,\mathbb C)=2n^2.
\]


% ============================================================
% SL
% ============================================================

\section{Specjalna grupa liniowa i algebra
\texorpdfstring{$\mathfrak{sl}(n,\mathbb K)$}{sl(n,K)}}

\index{grupa Liego!SL@$SL(n,\mathbb K)$}
\index{specjalna grupa liniowa}
\index{algebra Liego!sl@$\mathfrak{sl}(n,\mathbb K)$, specjalna liniowa}
\index{ślad macierzy}

\begin{definicja}
\emph{Specjalną grupą liniową} nazywamy
\[
SL(n,\mathbb K)
=
\{A\in GL(n,\mathbb K):\det A=1\}.
\]
Jej algebra Liego jest równa
\[
\mathfrak{sl}(n,\mathbb K)
=
\{X\in M_n(\mathbb K):\operatorname{tr}X=0\}.
\]
\end{definicja}

Warunek
\[
\operatorname{tr}X=0
\]
jest jednym niezależnym warunkiem liniowym. Dlatego
\[
\dim_{\mathbb R}SL(n,\mathbb R)
=
\dim_{\mathbb R}\mathfrak{sl}(n,\mathbb R)
=
n^2-1,
\]
a w przypadku zespolonym
\[
\dim_{\mathbb C}SL(n,\mathbb C)
=
\dim_{\mathbb C}\mathfrak{sl}(n,\mathbb C)
=
n^2-1.
\]
W szczególności
\[
\dim_{\mathbb R}SL(n,\mathbb C)=2(n^2-1).
\]


% ============================================================
% O i SO
% ============================================================

\section{Grupy ortogonalne i algebra
\texorpdfstring{$\mathfrak{so}(n)$}{so(n)}}
\label{sec:app-grupa-o}

\index{grupa Liego!O@$O(n)$}
\index{grupa Liego!SO@$SO(n)$}
\index{grupa ortogonalna}
\index{specjalna grupa ortogonalna}
\index{algebra Liego!so@$\mathfrak{so}(n)$, ortogonalna}
\index{macierz!antysymetryczna}

\begin{definicja}
\emph{Grupą ortogonalną} nazywamy
\[
O(n)
=
\{A\in GL(n,\mathbb R):A^TA=I\},
\]
natomiast \emph{specjalną grupą ortogonalną}
\[
SO(n)
=
\{A\in O(n):\det A=1\}.
\]
\end{definicja}

Obie grupy mają tę samą algebrę Liego:
\[
\operatorname{Lie}(O(n))
=
\operatorname{Lie}(SO(n))
=
\mathfrak{so}(n),
\]
gdzie
\[
\mathfrak{so}(n)
=
\{X\in M_n(\mathbb R):X^T=-X\}.
\]

Macierz antysymetryczna ma zerową przekątną, a dla $i<j$
współczynnik $x_{ij}$ może być wybrany dowolnie, po czym
\[
x_{ji}=-x_{ij}.
\]
Liczba niezależnych parametrów wynosi więc
\[
\binom n2=\frac{n(n-1)}2.
\]
Stąd
\[
\boxed{
\dim O(n)
=
\dim SO(n)
=
\dim\mathfrak{so}(n)
=
\frac{n(n-1)}2.
}
\]

Grupa $O(n)$ ma dwie składowe spójności, odpowiadające warunkom
$\det A=1$ i $\det A=-1$. Algebra Liego nie wykrywa tej różnicy,
ponieważ opisuje lokalną strukturę grupy jedynie w otoczeniu
elementu neutralnego.


% ============================================================
% SO(n,C)
% ============================================================

\section{Zespolona grupa ortogonalna}

\index{grupa Liego!SOC@$SO(n,\mathbb C)$}
\index{algebra Liego!soC@$\mathfrak{so}(n,\mathbb C)$, ortogonalna zespolona}

Definiujemy
\[
SO(n,\mathbb C)
=
\{A\in GL(n,\mathbb C):
A^TA=I,\ \det A=1\}.
\]
Jej algebra Liego ma postać
\[
\mathfrak{so}(n,\mathbb C)
=
\{X\in M_n(\mathbb C):X^T=-X\}.
\]

Stąd
\[
\boxed{
\dim_{\mathbb C}SO(n,\mathbb C)
=
\dim_{\mathbb C}\mathfrak{so}(n,\mathbb C)
=
\frac{n(n-1)}2.
}
\]
Jako rzeczywista grupa Liego ma ona wymiar
\[
n(n-1).
\]


% ============================================================
% U i SU
% ============================================================

\section{Grupy unitarne i ich algebry Liego}
\label{sec:app-grupa-u}

\index{grupa Liego!U@$U(n)$}
\index{grupa Liego!SU@$SU(n)$}
\index{grupa unitarna}
\index{specjalna grupa unitarna}
\index{algebra Liego!u@$\mathfrak u(n)$, unitarna}
\index{algebra Liego!su@$\mathfrak{su}(n)$, specjalna unitarna}
\index{macierz!antyhermitowska}

Przez
\[
A^*:=\overline A^{\,T}
\]
oznaczamy sprzężenie hermitowskie macierzy.

\begin{definicja}
\emph{Grupą unitarną} nazywamy
\[
U(n)
=
\{A\in GL(n,\mathbb C):A^*A=I\},
\]
natomiast \emph{specjalną grupą unitarną}
\[
SU(n)
=
\{A\in U(n):\det A=1\}.
\]
\end{definicja}

Ich algebry Liego są odpowiednio równe
\[
\mathfrak u(n)
=
\{X\in M_n(\mathbb C):X^*=-X\}
\]
oraz
\[
\mathfrak{su}(n)
=
\{X\in M_n(\mathbb C):
X^*=-X,\ \operatorname{tr}X=0\}.
\]

Macierz antyhermitowska ma czysto urojone elementy diagonalne,
co daje $n$ parametrów rzeczywistych. Dla każdej pary $i<j$
element $x_{ij}\in\mathbb C$ jest dowolny, a
\[
x_{ji}=-\overline{x_{ij}}.
\]
Każda taka para daje dwa parametry rzeczywiste. Zatem
\[
n+2\binom n2=n^2.
\]
Otrzymujemy
\[
\boxed{
\dim_{\mathbb R}U(n)
=
\dim_{\mathbb R}\mathfrak u(n)
=
n^2.
}
\]

Warunek $\operatorname{tr}X=0$ usuwa jeden rzeczywisty parametr,
dlatego
\[
\boxed{
\dim_{\mathbb R}SU(n)
=
\dim_{\mathbb R}\mathfrak{su}(n)
=
n^2-1.
}
\]

W szczególności
\[
\dim SU(2)=3,
\qquad
\dim SU(3)=8.
\]

Algebry $\mathfrak u(n)$ i $\mathfrak{su}(n)$ są naturalnie
rzeczywistymi algebrami Liego, mimo że składają się z macierzy
zespolonych.


% ============================================================
% Sp(2n,R)
% ============================================================

\section{Rzeczywista grupa symplektyczna}

\index{grupa Liego!SpR@$Sp(2n,\mathbb R)$}
\index{grupa symplektyczna!rzeczywista}
\index{algebra Liego!spR@$\mathfrak{sp}(2n,\mathbb R)$, symplektyczna rzeczywista}
\index{macierz!symplektyczna}

Niech
\[
J=
\begin{pmatrix}
0&I_n\\
-I_n&0
\end{pmatrix}.
\]

\begin{definicja}
\emph{Rzeczywistą grupą symplektyczną} nazywamy
\[
Sp(2n,\mathbb R)
=
\{A\in GL(2n,\mathbb R):A^TJA=J\}.
\]
Jej algebra Liego jest równa
\[
\mathfrak{sp}(2n,\mathbb R)
=
\{X\in M_{2n}(\mathbb R):X^TJ+JX=0\}.
\]
\end{definicja}

Jeśli
\[
X=
\begin{pmatrix}
A&B\\
C&D
\end{pmatrix},
\]
to warunek
\[
X^TJ+JX=0
\]
jest równoważny
\[
D=-A^T,
\qquad
B=B^T,
\qquad
C=C^T.
\]
Każdy element ma zatem postać
\[
X=
\begin{pmatrix}
A&B\\
C&-A^T
\end{pmatrix},
\]
gdzie $A$ jest dowolną macierzą $n\times n$, a $B$ i $C$ są
symetryczne.

Wobec tego
\[
\dim\mathfrak{sp}(2n,\mathbb R)
=
n^2
+
2\frac{n(n+1)}2
=
n(2n+1).
\]
Stąd
\[
\boxed{
\dim_{\mathbb R}Sp(2n,\mathbb R)
=
\dim_{\mathbb R}\mathfrak{sp}(2n,\mathbb R)
=
n(2n+1).
}
\]


% ============================================================
% Sp(2n,C)
% ============================================================

\section{Zespolona grupa symplektyczna}

\index{grupa Liego!SpC@$Sp(2n,\mathbb C)$}
\index{grupa symplektyczna!zespolona}
\index{algebra Liego!spC@$\mathfrak{sp}(2n,\mathbb C)$, symplektyczna zespolona}

Definiujemy
\[
Sp(2n,\mathbb C)
=
\{A\in GL(2n,\mathbb C):A^TJA=J\},
\]
oraz
\[
\mathfrak{sp}(2n,\mathbb C)
=
\{X\in M_{2n}(\mathbb C):X^TJ+JX=0\}.
\]

Ten sam rachunek blokowy daje
\[
\boxed{
\dim_{\mathbb C}Sp(2n,\mathbb C)
=
\dim_{\mathbb C}\mathfrak{sp}(2n,\mathbb C)
=
n(2n+1).
}
\]
Jako rzeczywista grupa Liego ma ona wymiar
\[
2n(2n+1).
\]


% ============================================================
% compact Sp(n)
% ============================================================

\section{Zwarta grupa symplektyczna}

\index{grupa Liego!Sp@$Sp(n)$}
\index{grupa symplektyczna!zwarta}
\index{algebra Liego!sp@$\mathfrak{sp}(n)$, symplektyczna zwarta}
\index{kwaterniony!grupa symplektyczna}

\begin{definicja}
\emph{Zwartą grupą symplektyczną} nazywamy
\[
Sp(n)
=
\{A\in M_n(\mathbb H):A^*A=I\},
\]
gdzie $\mathbb H$ oznacza algebrę kwaternionów.
Jej algebra Liego ma postać
\[
\mathfrak{sp}(n)
=
\{X\in M_n(\mathbb H):X^*=-X\}.
\]
\end{definicja}

Równoważnie $\mathfrak{sp}(n)$ można zrealizować jako podalgebrę
$M_{2n}(\mathbb C)$:
\[
\mathfrak{sp}(n)
=
\left\{
X\in M_{2n}(\mathbb C):
X^*=-X,\quad X^TJ+JX=0
\right\}.
\]

Jej wymiar rzeczywisty wynosi
\[
\boxed{
\dim_{\mathbb R}Sp(n)
=
\dim_{\mathbb R}\mathfrak{sp}(n)
=
n(2n+1).
}
\]

Należy odróżniać
\[
Sp(n),
\qquad
Sp(2n,\mathbb R),
\qquad
Sp(2n,\mathbb C).
\]
Pierwsze dwie są rzeczywistymi grupami Liego tego samego wymiaru,
ale mają inną geometrię globalną i odpowiadają różnym rzeczywistym
formom zespolonej algebry
\[
\mathfrak{sp}(2n,\mathbb C).
\]


% ============================================================
% PGL
% ============================================================

\section{Projektywna grupa liniowa}
\label{sec:app-grupa-pgl}

\index{grupa Liego!PGL@$PGL(n,\mathbb K)$}
\index{projektywna grupa liniowa}
\index{algebra Liego!pgl@$\mathfrak{pgl}(n,\mathbb K)$}

\begin{definicja}
\emph{Projektywną grupą liniową} nazywamy
\[
PGL(n,\mathbb K)
=
GL(n,\mathbb K)/Z(GL(n,\mathbb K)),
\]
gdzie
\[
Z(GL(n,\mathbb K))
=
\{\lambda I:\lambda\in\mathbb K_*\}.
\]
\end{definicja}

Ponieważ grupa niezerowych skalarów ma wymiar $1$ nad $\mathbb K$,
otrzymujemy
\[
\dim_{\mathbb R}PGL(n,\mathbb R)=n^2-1,
\]
oraz
\[
\dim_{\mathbb C}PGL(n,\mathbb C)=n^2-1.
\]

Jej algebra Liego ma postać
\[
\mathfrak{pgl}(n,\mathbb K)
=
\mathfrak{gl}(n,\mathbb K)/\mathbb K I.
\]
Dla $\mathbb K=\mathbb R$ lub $\mathbb C$ zachodzi izomorfizm
\[
\mathfrak{pgl}(n,\mathbb K)
\simeq
\mathfrak{sl}(n,\mathbb K).
\]
Wysyła on klasę $X+\mathbb K I$ na macierz
$X-(\operatorname{tr}X/n)I$: dodanie skalara
do $X$ nie zmienia wyniku, a każda macierz
o śladzie zerowym jest osiągana.


% ============================================================
% PSL
% ============================================================

\section{Projektywna specjalna grupa liniowa}

\index{grupa Liego!PSL@$PSL(n,\mathbb K)$}
\index{projektywna specjalna grupa liniowa}

\begin{definicja}
\emph{Projektywną specjalną grupą liniową} nazywamy
\[
PSL(n,\mathbb K)
=
SL(n,\mathbb K)/Z(SL(n,\mathbb K)).
\]
\end{definicja}

Centrum $SL(n,\mathbb K)$ składa się z macierzy
$\lambda I$ spełniających $\lambda^n=1$.
Jest więc skończone, dlatego przejście do
ilorazu nie zmienia wymiaru:
\[
\dim_{\mathbb R}PSL(n,\mathbb R)=n^2-1,
\]
oraz
\[
\dim_{\mathbb C}PSL(n,\mathbb C)=n^2-1.
\]

Jej algebra Liego jest równa
\[
\mathfrak{sl}(n,\mathbb K).
\]


% ============================================================
% PU
% ============================================================

\section{Projektywna grupa unitarna}

\index{grupa Liego!PU@$PU(n)$}
\index{projektywna grupa unitarna}

\begin{definicja}
\emph{Projektywną grupą unitarną} nazywamy
\[
PU(n)
=
U(n)/U(1),
\]
gdzie
\[
U(1)=\{\lambda I:|\lambda|=1\}
\]
jest centrum $U(n)$.
\end{definicja}

Ponieważ
\[
\dim_{\mathbb R}U(1)=1,
\]
otrzymujemy
\[
\boxed{
\dim_{\mathbb R}PU(n)=n^2-1.
}
\]

Jej algebra Liego jest izomorficzna z
\[
\mathfrak{su}(n).
\]


% ============================================================
% PSp compact
% ============================================================

\section{Projektywna zwarta grupa symplektyczna}

\index{grupa Liego!PSp@$PSp(n)$}
\index{projektywna grupa symplektyczna}

Centrum zwartej grupy $Sp(n)$ jest równe
\[
Z(Sp(n))=\{\pm I\}.
\]

\begin{definicja}
Definiujemy
\[
PSp(n)
=
Sp(n)/\{\pm I\}.
\]
\end{definicja}

Ponieważ $\{\pm I\}$ jest grupą dyskretną,
iloraz nie zmienia wymiaru:
\[
\boxed{
\dim_{\mathbb R}PSp(n)
=
n(2n+1).
}
\]

Jej algebra Liego jest równa
\[
\mathfrak{sp}(n).
\]


% ============================================================
% GSp
% ============================================================

\section{Grupa symplektycznych podobieństw}

\index{grupa Liego!GSp@$GSp(2n,\mathbb K)$}
\index{grupa symplektycznych podobieństw}

Standardowym oznaczeniem grupy zachowującej formę symplektyczną
z dokładnością do niezerowego skalara jest $GSp$, a nie $SPGL$.

\begin{definicja}
\emph{Grupą symplektycznych podobieństw} nazywamy
\[
GSp(2n,\mathbb K)
=
\left\{
A\in GL(2n,\mathbb K):
A^TJA=\lambda J
\text{ dla pewnego }
\lambda\in\mathbb K_*
\right\}.
\]
\end{definicja}

Skalar $\lambda$ jest jednoznacznie wyznaczony przez $A$ i nazywa
się \emph{mnożnikiem symplektycznym}.

Mnożnik daje surjektywny homomorfizm
$GSp(2n,\mathbb K)\to\mathbb K_*$ z jądrem
$Sp(2n,\mathbb K)$: wartość $\lambda$ przyjmuje
na macierzy $\operatorname{diag}(I_n,\lambda I_n)$.
Iloraz ma jeden parametr nad $\mathbb K$, dlatego
\[
\boxed{
\dim_{\mathbb K}GSp(2n,\mathbb K)
=
n(2n+1)+1.
}
\]

W szczególności
\[
\dim_{\mathbb R}GSp(2n,\mathbb R)
=
n(2n+1)+1,
\]
oraz
\[
\dim_{\mathbb C}GSp(2n,\mathbb C)
=
n(2n+1)+1.
\]

Jej algebra Liego składa się z macierzy spełniających
\[
X^TJ+JX=\mu J
\]
dla pewnego $\mu\in\mathbb K$.


% ============================================================
% PGSp
% ============================================================

\section{Projektywna grupa symplektycznych podobieństw}

\index{grupa Liego!PGSp@$PGSp(2n,\mathbb K)$}

Definiujemy
\[
PGSp(2n,\mathbb K)
=
GSp(2n,\mathbb K)/(\mathbb K_* I).
\]

Ilorazujemy przez jednowymiarową centralną grupę macierzy skalarnych
\[
\{\lambda I:\lambda\in\mathbb K_*\}.
\]
W konsekwencji
\[
\boxed{
\dim_{\mathbb K}PGSp(2n,\mathbb K)
=
n(2n+1).
}
\]


% ============================================================
% TABELA ZBIORCZA
% ============================================================

\section{Tabela zbiorcza}

\[
\begin{array}{c|c|c}
\text{grupa}
&
\text{algebra Liego}
&
\text{wymiar}
\\
\hline

GL(n,\mathbb R)
&
\mathfrak{gl}(n,\mathbb R)
&
n^2
\\[1mm]

GL(n,\mathbb C)
&
\mathfrak{gl}(n,\mathbb C)
&
n^2\ \text{nad }\mathbb C
\\[1mm]

SL(n,\mathbb R)
&
\mathfrak{sl}(n,\mathbb R)
&
n^2-1
\\[1mm]

SL(n,\mathbb C)
&
\mathfrak{sl}(n,\mathbb C)
&
n^2-1\ \text{nad }\mathbb C
\\[1mm]

O(n)
&
\mathfrak{so}(n)
&
\dfrac{n(n-1)}2
\\[3mm]

SO(n)
&
\mathfrak{so}(n)
&
\dfrac{n(n-1)}2
\\[3mm]

SO(n,\mathbb C)
&
\mathfrak{so}(n,\mathbb C)
&
\dfrac{n(n-1)}2\ \text{nad }\mathbb C
\\[3mm]

U(n)
&
\mathfrak u(n)
&
n^2
\\[1mm]

SU(n)
&
\mathfrak{su}(n)
&
n^2-1
\\[1mm]

Sp(2n,\mathbb R)
&
\mathfrak{sp}(2n,\mathbb R)
&
n(2n+1)
\\[1mm]

Sp(2n,\mathbb C)
&
\mathfrak{sp}(2n,\mathbb C)
&
n(2n+1)\ \text{nad }\mathbb C
\\[1mm]

Sp(n)
&
\mathfrak{sp}(n)
&
n(2n+1)
\\[1mm]

PGL(n,\mathbb R)
&
\mathfrak{pgl}(n,\mathbb R)
\simeq
\mathfrak{sl}(n,\mathbb R)
&
n^2-1
\\[1mm]

PGL(n,\mathbb C)
&
\mathfrak{pgl}(n,\mathbb C)
\simeq
\mathfrak{sl}(n,\mathbb C)
&
n^2-1\ \text{nad }\mathbb C
\\[1mm]

PSL(n,\mathbb K)
&
\mathfrak{sl}(n,\mathbb K)
&
n^2-1
\\[1mm]

PU(n)
&
\mathfrak{su}(n)
&
n^2-1
\\[1mm]

PSp(n)
&
\mathfrak{sp}(n)
&
n(2n+1)
\\[1mm]

GSp(2n,\mathbb K)
&
\mathfrak{gsp}(2n,\mathbb K)
&
n(2n+1)+1
\\[1mm]

PGSp(2n,\mathbb K)
&
\mathfrak{pgsp}(2n,\mathbb K)
&
n(2n+1)
\end{array}
\]

\paragraph{Przykłady w małym wymiarze.}
Dla $n=1$ mamy $GL(1,\mathbb K)=\mathbb K_*$,
$SL(1,\mathbb K)=\{1\}$, $U(1)=S^1$ i $SU(1)=\{1\}$.
Grupa $O(2)$ składa się z obrotów i odbić płaszczyzny,
a $SO(2)$ z samych obrotów; jej wymiar $1$ odpowiada
kątowi obrotu. Z kolei $SU(2)\cong S^3$ obliczyliśmy
w Przykładzie~\ref{prz:su2-liego}.

W wymiarze symplektycznym $2$ zachodzą
$Sp(2,\mathbb K)=SL(2,\mathbb K)$ i
$GSp(2,\mathbb K)=GL(2,\mathbb K)$,
bo $A^TJA=(\det A)J$ dla każdej macierzy $2\times2$.
Stąd $PGSp(2,\mathbb K)=PGL(2,\mathbb K)$;
obie grupy projektowe działają na prostej
$\mathbb P^1(\mathbb K)$ przez klasy macierzy.
W tej samej akcji $PSL(2,\mathbb K)$ jest obrazem
podgrupy $SL(2,\mathbb K)$.
Zwarta $Sp(1)$ to grupa jednostkowych kwaternionów,
izomorficzna z $SU(2)$, a $PSp(1)$ jest ilorazem
przez $\{\pm1\}$, izomorficznym z $SO(3)$.
Podobnie $PU(1)$ jest grupą trywialną.
Te przykłady pokazują, dlaczego odmienne oznaczenia
$Sp(2n,\mathbb K)$ i $Sp(n)$ trzeba rozróżniać.


\begin{uwaga}
Iloraz grupy Liego przez podgrupę dyskretną nie zmienia wymiaru.
Dlatego na przykład
\[
\dim PSL(n,\mathbb K)
=
\dim SL(n,\mathbb K)
\]
oraz
\[
\dim PSp(n)
=
\dim Sp(n).
\]
Natomiast w przypadku
\[
PGL(n,\mathbb K)
=
GL(n,\mathbb K)/\mathbb K_*
\]
ilorazujemy przez jednowymiarową grupę Liego, dlatego
\[
\dim_{\mathbb K}PGL(n,\mathbb K)
=
n^2-1.
\]
\end{uwaga}

\begin{uwaga}
Należy uważać na oznaczenia związane z grupami symplektycznymi.
W literaturze występują równolegle:
\[
Sp(2n,\mathbb R),
\qquad
Sp(2n,\mathbb C),
\qquad
Sp(n).
\]
Pierwsza jest rzeczywistą grupą symplektyczną, druga jej zespolonym
odpowiednikiem, natomiast $Sp(n)$ oznacza zwartą grupę
kwaternionowo-unitarną. Wszystkie mają algebry Liego o wymiarze
$n(2n+1)$ nad odpowiednim ciałem, lecz nie są tym samym obiektem.
\end{uwaga}
